\chapter[chap:phase4]{Phase 4: Akteure und Zusammenarbeit}

\section{Ziel}

Schreibhoheit, Pull Requests und Freigaben sollen im System erzwungen werden, nicht nur
empfohlen.
Dieses Kapitel setzt ADR-002 und ADR-004 um und stützt sich für das Merge-Gate auf ADR-019.

\section{Umfang}

\startitemize
\item Lease pro Branch, die jede schreibende Operation erneuert, mit Ablaufzeit.
\item Owner-Override nur durch einen Menschen mit der Rolle Owner, signiert, mit Begründung und
Eintrag im Operationslog, ohne Umgehung des Merge-Gates.
\item Pull Request als Folge von Einträgen im Operationslog, mit Reviewern, Prüfungen und
Merge-Regeln.
\item Ein Merge-Gate ohne globale Sperre (ADR-019).
\item Freigaben mit Akteur, Modellversion, geprüftem Commit und Ergebnis, signiert.
\item Verbot der Selbstfreigabe.
\item Risikostufen 0 bis 3, festgelegt in einer Policy-Datei, die selbst Stufe 3 ist.
\item Automatische Eskalation bei Lockerung bestehender Tests.
\stopitemize

\section{Schreibhoheit und Pull Request}

Dieser Abschnitt beschreibt ADR-004, das für Phase 4 angenommen ist.
Umgesetzt ist davon noch nichts.
Zwei Regeln des ADR bestimmen den Weg einer Änderung auf \type{main}.
Jeder Branch hat genau einen Schreibinhaber, und \type{main} ändert sich nur durch einen Pull
Request.

{\bf Die Lease.}
Der Kern erzwingt die Schreibhoheit mit einer Lease, die jede schreibende Operation erneuert und
die nach einer Ablaufzeit endet (\in{Abbildung}[fig:branch-lease]).

\startitemize[n]
\item Ein Akteur erwirbt die Lease eines Branches und wird sein Schreibinhaber.
\item Solange er schreibt, verlängert jede Operation die Lease.
\item Stürzt er ab, bleibt die Verlängerung aus.
\item Die Lease läuft ab, und der Branch ist frei.
\item Ein anderer Akteur erwirbt sie.
\stopitemize

Schreibt ein zweiter Akteur, solange die Lease gilt, lehnt der Kern das ab.
Wer nicht auf den Ablauf warten will, braucht einen Owner-Override.
Ihn darf nur ein Mensch mit der Rolle Owner auslösen, signiert und mit einer Begründung.
Er hinterlässt einen Eintrag im Operationslog, den der Purge nie anfasst,
weil er keinen Commit nennt, und er umgeht das Merge-Gate nie.
Die Lease dauert standardmässig 60 Minuten; die Policy kann sie ändern, innerhalb fester Grenzen
von 5 Minuten bis 24 Stunden.
Nach dem Ablauf darf jeder mit Schreibrecht sie übernehmen, auch nach dem Absturz eines Agenten.
Im offenen Modus und auf geschützten Branches gibt es keine Lease;
sie entsteht mit der ersten schreibenden Operation, sobald das Register mehr als einen Akteur nennt.

\startplacefigure[title={Zeitachse einer Branch-Lease},
    reference=fig:branch-lease,location=here]
\externalfigure[images/branch-lease.pdf][width=\textwidth]
\stopplacefigure

{\bf Der Pull Request.}
Der Schreibinhaber eröffnet den Pull Request (\in{Abbildung}[fig:pull-request-flow]).
Die Risikostufe folgt aus allen Pfaden, die die Änderung berührt, und der höchste Wert gilt.
Sie bestimmt, wer freigeben muss, und kein Akteur gibt seine eigene Änderung frei.
Ist die Freigabe vollständig, wendet das Merge-Gate die Regel an (\in{Abschnitt}[sec:merge-gate])
und bewegt \type{main}.
Ein Konflikt schickt die Änderung an ihren Autor zurück.
Hat sich \type{main} inzwischen bewegt, wiederholt das Gate den Merge und die Regel
auf der neuen Spitze.

\startplacefigure[title={Ablauf eines Pull Requests bis zur Zeigerbewegung von main},
    reference=fig:pull-request-flow,location=here]
\externalfigure[images/pull-request-flow.pdf][width=\textwidth]
\stopplacefigure

\section[sec:merge-gate]{Das Merge-Gate ohne globale Sperre}

Eine Merge-Queue mit globaler Sperre hält jede Änderung an, solange eine andere getestet wird.
Bei vielen Schreibern, Menschen wie Agenten, wächst die Wartezeit mit ihrer Zahl und mit der
Dauer einer Prüfung.
Die Frage ist, was eine grüne Prüfung noch heisst, wenn während ihres Laufs viele andere
Änderungen landen.
ADR-019 beantwortet sie; er ist angenommen, aber noch nicht umgesetzt.

{\bf Der Ablauf.}

\startitemize[n]
\item Das Gate liest die Spitze \type{T} von \type{main} und führt die Änderung darauf zusammen.
Das Ergebnis ist der Commit \type{M}, der Squash-Commit der Änderung auf \type{T}
(ADR-006); \type{main} bleibt linear.
\item Es wendet die Regel nach Risikostufe (\in{Tabelle}[tab:gate-rule])
auf die vorhandenen Prüfergebnisse an.
\item Es setzt \type{main} von \type{T} auf \type{M}, per Compare-and-swap (ADR-008).
Hat eine andere Änderung \type{main} inzwischen bewegt, scheitert die Zeigerbewegung (\type{Stale},
oder \type{Locked}, solange deren Sperre noch gehalten wird) und das Gate beginnt
bei Schritt 1 mit der neuen Spitze.
\stopitemize

Während eine Prüfung läuft, hält niemand eine Sperre.
Die einzige Serialisierung ist die Zeigerbewegung von \type{main};
die Sperre der Referenzen gilt für eine Änderung, nie für eine Prüfung.
Das Gate schreibt die Objekte des Merges vor der Sperre.
Unter der Sperre bleiben die Prüfung der Spitze, der Eintrag des Operation Logs und die
Zeigerbewegung (ADR-008, ADR-011).
Weil ein Purge unreferenzierte Objekte löschen könnte, prüft das Gate unter der Sperre,
ob \type{M} im Speicher liegt, und beginnt sonst bei Schritt 1 von vorn.
\type{track}, \type{commit} und \type{op restore} nehmen die Sperre weiterhin vor ihrem ersten
Objekt.

{\bf Was ein Prüfergebnis aussagt.}
Eine Prüfung ist eine Aussage über einen Commit: Sein Tree hat bestanden oder nicht.
Über den Commit, zu dem \type{main} wird, wenn die Änderung landet, sagt sie nichts, denn andere
Änderungen können vorher landen.
Ein Ergebnis hält deshalb zwei Commits fest, den geprüften und die Spitze von \type{main}, auf die
er für die Prüfung gelegt wurde.
Ob es nach einer Bewegung von \type{main} genügt, bestimmt die Risikostufe aus ADR-004:

\startplacetable[title={Gültigkeit eines Prüfergebnisses nach einer Bewegung von main},
reference=tab:gate-rule]
\startxtable[frame=off,option=stretch]
\startxrow[topframe=on,bottomframe=on]
\startxcell {\bf Stufe} \stopxcell
\startxcell {\bf Regel, wenn \type{main} sich nach der Prüfung bewegt hat} \stopxcell
\stopxrow
\startxrow
\startxcell 0, 1 \stopxcell
\startxcell Das Ergebnis bleibt gültig, wenn der Merge auf die neue Spitze sauber ist.
\stopxcell
\stopxrow
\startxrow[bottomframe=on]
\startxcell 2, 3 \stopxcell
\startxcell Das Ergebnis gilt nur für die Spitze, gegen die es lief; nach einer Bewegung laufen
die Prüfungen auf dem neuen Merge noch einmal. \stopxcell
\stopxrow
\stopxtable
\stopplacetable

Sauber heisst: Der Tree-Merge liefert genau einen Tree (\in{Abschnitt}[sec:tree-merge]).

{\bf Konflikte.}
Ein Konflikt gelangt nie nach \type{main}.
Hat der Merge in Schritt 1 einen, lehnt das Gate ab und gibt die Änderung an ihren Autor zurück.
Der Autor, der Schreibinhaber des Branches (ADR-004), spielt ihn mit \type{rebase} auf die neue
Spitze neu auf; die umgeschriebenen Commits tragen den Konflikt als Daten
(\in{Kapitel}[chap:phase3]).
Er löst ihn und pusht erneut.
Lehnt \type{rebase} ab, etwa bei einem Merge-Commit im Bereich, ist das ein gewöhnlicher Fehler,
und der Autor entscheidet von Hand.
Das Konfliktmodell der Phase 3 ist die Voraussetzung: Mit Markierungen in den Dateien liesse sich
ein Konflikt nicht als Commit zurückgeben.
\in{Abbildung}[fig:conflict-gate] zeigt den Konflikt im Merge und den semantischen Konflikt,
die das Gate an verschiedenen Stellen erkennt.

\startplacefigure[title={Konflikte am Merge-Gate: im Merge und auf der Spitze von main},
    reference=fig:conflict-gate,location=here]
\externalfigure[images/conflict-gate.pdf][width=\textwidth]
\stopplacefigure

{\bf Die Prüfung nach dem Merge.}
Zwei Änderungen können sauber zusammengehen und sich doch gegenseitig brechen, ein semantischer
Konflikt.
Nur ein Lauf auf der Spitze von \type{main} erkennt ihn.
Der Lauf deckt die Spitze zum Zeitpunkt seines Starts ab; was währenddessen landet, deckt der
nächste Lauf ab.
Ist die Spitze rot, behebt eine neue Änderung den Fehler, ein Fix oder ein Revert, und nicht ein
Umschreiben von \type{main}.
Solange die letzte geprüfte Spitze rot ist, landen keine Änderungen der Stufen 2 und 3; die
Stufen 0 und 1 gehen weiter.
Ausgenommen von der Pause ist die Änderung, die die rote Spitze behebt, etwa der Revert des
Commits, den die Bisektion findet.
Sie wird an einer Erklärung erkannt, nicht an ihrem Inhalt: Das \type{pr-open} nennt den roten
Commit (\type{repairs <commit>}).
Das Gate prüft, dass dieser Commit zwischen der letzten grünen und der ersten roten Spitze liegt
und dass die Prüfungen der Behebung mit der aktuellen Spitze als Basis bestanden haben.
Die Ursache in einem Fenster ungeprüfter Commits findet eine Bisektion über den linearen Verlauf.

{\bf Zwei Grenzen.}
Beide Werte stehen in der Policy-Datei; sie zu ändern ist Stufe 3, also ein Mensch mit der Rolle
Owner und ein ADR.
Ein Agent kann die Grenzen, unter denen seine eigene Änderung landet, nicht lockern.

\startitemize
\item {\bf Reservierung.}
Nach \type{N} ungültig gewordenen Läufen in Folge (Standard 3) erhält eine Änderung der Stufe 2
oder 3 eine Reservierung.
Merges der Stufen 0 und 1 melden dann \type{Paused} und warten,
bis die Änderung gelandet, gescheitert oder geschlossen ist.
Die Reservierung endet nach der doppelten Dauer des letzten Prüflaufs,
damit eine hängende Änderung \type{main} nicht blockiert.
Es gibt höchstens eine, und die älteste Änderung gewinnt.
\item {\bf Fenster ungeprüfter Commits.}
Die Policy kennt \type{max_unverified} (Standard 10).
Liegen so viele Commits nach der letzten geprüften Spitze, warten alle Merges der Stufen 0 bis 3,
bis ein Lauf auf der Spitze gestartet ist.
Die Behebung ist ausgenommen.
\stopitemize

{\bf Wer unabhängig ist.}
\quotation{Unabhängig} in ADR-004 heisst: verschiedene Akteur-IDs, zusammen mit dem Verbot der
Selbstfreigabe.
Zwei Agenten desselben Sponsors zählen als zwei Freigaben.
\type{model} und \type{session} beschreiben nur, daher lässt sich eine stärkere Unabhängigkeit
nicht prüfen; der Mensch bei den Stufen 2 und 3 bleibt der Anker.
Für die Stufen 2 und 3 zählt eine Prüfung (\type{check}) nur,
wenn ihr Schreiber weder die Änderung eröffnet
noch einen Commit in \type{base..head} verfasst hat
und kein Agent ist, dessen Sponsor einer von beiden ist;
bei den Stufen 0 und 1 zählt jeder registrierte Akteur.

\startinfobox
{\bf Was \quotation{main ist immer grün} jetzt heisst.}
Jeder Commit von \type{main} ist der saubere Merge einer Änderung, deren Prüfungen im Augenblick
der Zeigerbewegung nach der Regel oben bestanden hatten.
Die letzte geprüfte Spitze ist grün, oder eine Behebung läuft.
Zwischen dieser Spitze und der aktuellen liegen nur die Commits, die während eines Laufs landeten;
ihre Zahl ist durch \type{max_unverified} begrenzt, aber nicht null.
Grün sagt \quotation{dieser Commit hat bestanden}, nicht \quotation{der Stand von \type{main}
besteht jetzt}.
\stopinfobox

Der Entwurf ersetzt keine Queue durch eine andere.
Der Durchsatz hängt von der Rate der sauberen Zeigerbewegungen ab,
nicht von der Dauer der Prüfungen.
Dafür treten rote Spitzen auf, im Verhältnis zur Rate der semantischen Konflikte, und ein Teil der
Änderungen muss repariert werden, statt dass alle warten.
Änderungen der Stufen 2 und 3 holen bei vielen Schreibern die Spitze womöglich nie ein;
dafür gibt es die Reservierung (nächster Absatz).
Verworfen sind eine Queue mit globaler Sperre, eine Queue mit spekulativen Batches, die Prüfung
jeder Änderung bei jeder Bewegung von \type{main} und der Verzicht auf die Prüfung nach dem Merge;
die Gründe stehen in ADR-019.

\section[sec:phase4-open]{Offene Fragen}

\startitemize
\item Schnellweg für triviale Änderungen: Wer darf ihn nutzen?
\item Merge-Gate: Form und Ablage der Prüfergebnisse und die Frage, ob ein Ergebnis eine Menge
getesteter Basis-Commits nennen darf (offene Punkte in ADR-019).
Ein billigerer Test als der saubere Merge ist ausserhalb des Umfangs der ersten Version.
\stopitemize
